% DONO: todos (fechar juntos no fim)   ·   ESPAÇO: ~0,75 página
\section{Introdução}
% O enunciado pede: o domínio, o problema e por que uma linguagem própria
% faz sentido EM VEZ DE UMA BIBLIOTECA.
%
% Pontos a cobrir:
% - Promoções de loja virtual são regras de negócio que se combinam.
%   Combinações não previstas (descontos que se somam, fretes que se
%   contradizem) só aparecem quando um carrinho real as dispara.
% - Plataformas validam um carrinho por vez, em execução. CUIDADO: não
%   afirmem nada sobre Shopify/VTEX sem citar a documentação delas
%   (ex.: regras de combinação de descontos). Não digam que o total fica
%   negativo: as plataformas limitam o preço a zero.
% - Tese: com uma linguagem de condições FECHADA (só qtd, total e
%   constantes), o compilador analisa o conjunto de regras para QUALQUER
%   carrinho possível. Numa biblioteca Python, a condição seria uma função
%   arbitrária, e não dá para decidir se duas funções arbitrárias valem
%   juntas. Citar Mernik et al. \cite{mernik2005} sobre quando criar uma DSL.
% - Contribuições (lista curta): a linguagem; a verificação de conflitos
%   por faixas; a avaliação com N programas.
% - Uma frase sobre a organização do artigo.
